% Work-MAT | 2. rocnik | pripravovana hodina 44
% Návaznost: uživatelské připravy.zip (hodiny 01–20), autorská učebnice na GitHubu.
% Realisticky.cz: didaktická inspirace v posloupnosti a typech úloh; zadání jsou vlastní.
\documentclass[11pt,a4paper]{article}
\input{../../sablony/pripravy/priprava.sty}
\nastavPripravu{2. ročník}{Logaritmické rovnice a nerovnice}{44}{Smíšené logaritmické úlohy}
\begin{document}
\blok{Motivace a návaznost}
Dnešní souhrnná hodina nutí volit mezi převodem na exponenciální tvar, spojením logaritmů a monotónností.
\poznamka{Orientační nabídka pro 45--55 minut: krátké zopakování a motivace, nový výklad, tři společné vzorové úlohy, samostatná práce. Vyber podle tempa třídy; zbývající zadání slouží jako rezerva.}
\blok{Výklad}
Postup: (1) podmínky, (2) určení vztahu a metody, (3) úpravy, (4) průnik výsledku s podmínkami, (5) kontrola.
\blok{Společné příklady}
\begin{prikladReseniBox}[title={Příklad 1}]
Vyřeš, resp. urči: $\log_3(x+1)>\log_3(4-x)$.
\tcblower
\[\begin{aligned}-1<x<4 \\ x+1>4-x \\ \frac{3}{2}<x<4\end{aligned}\]
\end{prikladReseniBox}
\begin{prikladReseniBox}[title={Příklad 2}]
Vyřeš, resp. urči: $\log_2(x-1)+\log_2 2=3$.
\tcblower
\[\begin{aligned}x>1 \\ 2(x-1)=8 \\ x=5\end{aligned}\]
\end{prikladReseniBox}
\poznamka{Dej žákům možnost označit typ úlohy barevně ve vlastním sešitě; u všech řešení zapisuj interval.}
\blok{Samostatné procvičování}
\begin{zadaniReseniBox}[title={Samostatná práce 1}]
\begin{enumerate}[label=\alph*)]
\item $\log_{2}(x+1)>\log_{2}(4-x)$.
\item $\log_{3}(x+2)>\log_{3}(13-x)$.
\item $\log_{5}(x+3)>\log_{5}(131-x)$.
\item $\log_{2}(x+4)>\log_{2}(10-x)$.
\item $\log_{3}(x+1)>\log_{3}(11-x)$.
\item $\log_{5}(x+2)>\log_{5}(129-x)$.
\end{enumerate}
\tcblower
\begin{enumerate}[label=\alph*)]
\item $\frac{3}{2}<x<4$.
\item $\frac{11}{2}<x<13$.
\item $\frac{128}{2}<x<131$.
\item $\frac{6}{2}<x<10$.
\item $\frac{10}{2}<x<11$.
\item $\frac{127}{2}<x<129$.
\end{enumerate}
\end{zadaniReseniBox}
\blok{Dobrovolné rozšíření}
\begin{zadaniReseniBox}[title={Rozšíření 1}]
\begin{enumerate}[label=\alph*)]
\item $Vytvoř vlastní úlohu tohoto typu a se spolužákem porovnej podmínky a výsledek.$.
\end{enumerate}
\tcblower
\begin{enumerate}[label=\alph*)]
\item Zkontroluj definice a argumenty..
\end{enumerate}
\end{zadaniReseniBox}
\blok{Úkol na doma}
U každého logaritmu vypiš podmínky argumentu a proveď zkoušku. U rovnic vyznač případně nepřípustné kořeny, u nerovnic odpověz intervalem.
\begin{zadaniReseniBox}[title={Domácí úkol}]
\begin{enumerate}[label=\alph*)]
\item $\log_{2}(x+3)>\log_{2}(8-x)$.
\item $\log_{3}(x+4)>\log_{3}(17-x)$.
\item $\log_{5}(x+1)>\log_{5}(127-x)$.
\item $\log_{2}(x+2)>\log_{2}(6-x)$.
\item $\log_{3}(x+3)>\log_{3}(15-x)$.
\item $\log_{5}(x+4)>\log_{5}(133-x)$.
\item $\log_{2}(x+1)>\log_{2}(4-x)$.
\item $\log_{3}(x+2)>\log_{3}(13-x)$.
\end{enumerate}
\tcblower
\begin{enumerate}[label=\alph*)]
\item $\frac{5}{2}<x<8$.
\item $\frac{13}{2}<x<17$.
\item $\frac{126}{2}<x<127$.
\item $\frac{4}{2}<x<6$.
\item $\frac{12}{2}<x<15$.
\item $\frac{129}{2}<x<133$.
\item $\frac{3}{2}<x<4$.
\item $\frac{11}{2}<x<13$.
\end{enumerate}
\end{zadaniReseniBox}
\poznamka{Domácí úkol obsahuje 8 položek; předpokládaná domácí práce přibližně 15--30 minut. Po větších blocích jsou vědomě promíchané dříve probrané typy.}
\end{document}
